%
% board_ordering.tex
%
% A Z specification of the beads applet's board ordering: which of two boards
% the applet keeps, and which of two boards the display ends up showing.
%
% The slot half of this question is settled and shipped: BoardSlot keeps the
% board whose load began last, under one lock, and that rule makes the slot
% monotone in BoardOrder place.  The glass half is not settled.  A push is a
% round trip to luxd taken outside every lock, and two pushers may take theirs
% in either order, so the display can be left showing a board older than one it
% has already shown -- and older than one the slot has already refused.
%
% This model states both halves as one order and checks them together.  The
% committed design (Sections 5-6) makes each push read the slot and write the
% glass as one serialised step.  Three fidelity variants, each in its own file,
% remove one part of that and reach the exact bad state:
%
%   docs/board_ordering_current_buggy.tex     -- the code as it stands
%   docs/board_ordering_gate_only_buggy.tex   -- a monotone push counter alone
%   docs/board_ordering_reread_only_buggy.tex -- re-reading the slot alone
%
% Source of record:
%   src/punt_lux/applets/board_order.py    -- BoardOrder: the place a load takes
%   src/punt_lux/applets/board_slot.py     -- BoardSlot: store, newer_of, held
%   src/punt_lux/applets/held_board.py     -- HeldBoard: answer, refreshed
%   src/punt_lux/applets/no_board.py       -- NoBoard: answer, refreshed
%   src/punt_lux/applets/board_channel.py  -- BoardChannel: raised, send
%   src/punt_lux/applets/beads_service.py  -- BeadsService: prefetch,
%                                             acknowledge, service
%
% Type-check:  fuzz -t docs/board_ordering.tex
% Model-check: make prob SPEC=docs/board_ordering.tex
%   The five safety properties are reachability goals, not state invariants, so
%   that the fidelity variants can reach the bad states.  Each must report the
%   goal NOT found; the exact predicates are in Section 7.
%
% Formatting follows the fuzz house rules: \quad~ for continuation indent,
% schema lines under 80 columns, predicates broken at \land/\lor.
%

\documentclass[a4paper,10pt,fleqn]{article}
\usepackage[margin=1in]{geometry}
\usepackage{fuzz}
\usepackage[colorlinks=true,linkcolor=blue,citecolor=blue,urlcolor=blue]{hyperref}

\begin{document}

\title{Board Ordering in the Beads Applet: a Z Specification}
\author{Formal model of slot order and glass order}
\date{August 2026}
\maketitle

% ============================================================
\section{Introduction}
% ============================================================

The applet holds one board and shows one board, and those are two different
things. What it holds lives in \texttt{BoardSlot}; what it shows lives on the
display, at the far end of a round trip to \texttt{luxd}. Two loads can be in
flight at once --- the warm-up on one worker thread and a click on another ---
so both the slot and the display can be written by either of them, in either
order.

The slot half is settled. A load takes a \texttt{BoardOrder} place when it
\emph{begins}, the store keeps whichever of two boards holds the later place,
and the whole read-compare-write runs under one lock. That rule makes the slot
monotone in place: however the two threads interleave, the slot ends up holding
the board built from the newest issues.

The display half is not settled, and it is not settled in the same way. A push
carries a board across a socket, which takes long enough that it cannot be done
under the slot's lock, so the write to the display is ordered by nothing at all.
Two facts follow, and both are defects:

\begin{itemize}
  \item A click's answer captures the board the slot holds and pushes it after
    raising the frame. A refresh that stores \emph{and} pushes a newer board in
    that window is overwritten on the glass by the older capture. If the click
    then stands down --- because the refresh it would have run is the one
    already in flight --- nothing repairs the display. The user is left reading
    a board the applet has already replaced.
  \item A refresh pushes the board it loaded \emph{before} the slot has had a
    chance to refuse it. A load that began earlier but returned later therefore
    puts its board on the glass, and only afterwards learns that the slot kept a
    newer one. The screen shows a board the slot has discarded.
\end{itemize}

This is the fourth ordering defect found on this path in successive review
rounds, which is the signal to stop patching and model the order. What follows
is that model. It abstracts everything about a board except its place, because
place is the whole of what these defects turn on.

\paragraph{What is deliberately absent.}
Two mechanisms in the code do not appear here, each for a stated reason.
\texttt{SingleFlight} restricts refreshes to one at a time; the model omits it,
which admits \emph{more} interleavings than the implementation can produce. That
is the conservative direction: a design proved safe over the larger set is safe
over the smaller one. Standing down is therefore not an operation but a path ---
a worker that answers and never refreshes --- and it is the path that makes the
first defect permanent rather than transient. The frame's own lifetime (whether
a frame is up to be raised) is also absent: it decides whether a placeholder is
worth pushing, not which of two boards wins.

% ============================================================
\section{Basic Types}
% ============================================================

\begin{zed}
  [WORKER]
\end{zed}

\texttt{WORKER} is the set of threads that load boards and push them. Two
suffice: every defect here is a race between two writers, and a third adds
interleavings without adding shapes. \texttt{DEFAULT\_SETSIZE} bounds the
carrier.

\begin{zed}
  PLACE == 0 \upto 3
\end{zed}

A board's place in the order loads began, as \texttt{BoardOrder} numbers them.
Place $0$ is the state holding no board: \texttt{NoBoard}, and the content it
renders --- the ``Loading issues\ldots'' placeholder, or the red message naming
why \texttt{bd} could not be read. Places $1$ to $3$ are loads, numbered as they
begin. Three loads are enough to exhibit every defect (two contending boards and
one more to supersede them); the bound is a single token to change.

Comparing places compares the order two loads \emph{began} in, which is the rule
\texttt{BoardOrder.after} implements and the rule the model inherits: a board is
as old as the query that read it, not as old as the moment it came back.

% ============================================================
\section{State}
% ============================================================

The state is flat: one schema, nine components. Four are real, four are ghosts
that exist only so a property of the transition relation can be stated as a
property of a state, and one --- \texttt{staged} --- is both.

The state predicate carries only structural bounds, true of every design in this
family. The five safety properties are deliberately \emph{absent} from it. Were
they state invariants, ProB would treat an operation that broke one as
\emph{disabled} rather than as a violation, and the fidelity variants could no
longer reach the bad states they exist to reach. They are checked as
reachability goals instead (Section~\ref{sec:verify}).

\begin{schema}{Board}
  began : PLACE \\
  loading : WORKER \pfun PLACE \\
  held : PLACE \\
  shown : PLACE \\
  staged : WORKER \pfun PLACE \\
  everheld : PLACE \\
  evershown : PLACE \\
  offered : PLACE \\
  overtaken : \power PLACE
\where
  held \leq began \\
  shown \leq began \\
  everheld \leq began \\
  evershown \leq began \\
  offered \leq began
\end{schema}

\texttt{began} is the counter \texttt{BoardOrder} draws places from: the highest
place any load has taken. \texttt{loading} maps each worker running a load to
the place that load took when it began, and a worker leaves the domain when its
board is stored. \texttt{held} is the slot --- the place of the board
\texttt{BoardSlot} is holding, $0$ for \texttt{NoBoard}. \texttt{shown} is the
glass: the place of the board the display is rendering.

\texttt{staged} is the push in progress. A push is not one act but two: a
pusher settles what it is going to write, and some time later the write lands.
A worker in $\dom staged$ has settled its content and not yet landed it, and
$staged~w$ is the place of the board it is going to write. This split is the
whole point of the model --- collapse it and every ordering defect on this path
becomes unstateable.

The four ghosts record history. \texttt{everheld} is the highest place the slot
has ever held and \texttt{evershown} the highest the glass has ever shown, so
that ``never goes backwards'' can be written as an equation. \texttt{offered} is
the highest place any push has settled on, so that ``the glass ends up showing
the newest board anybody tried to put there'' can be written as an inequality.
\texttt{overtaken} collects the places the slot \emph{refused}: a board offered
to the store that lost the comparison and was never authoritative. That is a
narrower set than ``boards the slot no longer holds'' --- a board that was
authoritative and has since been superseded is merely stale, whereas a board the
slot refused was never the applet's answer at all, and the display showing one
is a different and worse fact.

% ============================================================
\section{Initialization}
% ============================================================

\begin{schema}{Init}
  Board'
\where
  began' = 0 \\
  loading' = \emptyset \\
  held' = 0 \\
  shown' = 0 \\
  staged' = \emptyset \\
  everheld' = 0 \\
  evershown' = 0 \\
  offered' = 0 \\
  overtaken' = \emptyset
\end{schema}

A single initial state: no load has begun, the slot holds no board, the glass
shows none, and nothing is in flight. This is the session at the moment the
applet registers, before the warm-up has run.

% ============================================================
\section{Loading and Storing}
% ============================================================

\texttt{TakePlace} is a load beginning: \texttt{BoardOrder.beginning()} taking
the next place from the counter, before the query runs. It is one operation for
both loaders --- the warm-up in \texttt{BeadsService.prefetch} and the refresh
behind a click --- because the place is taken the same way by each, and which of
them took it is not what decides the order.

\begin{schema}{TakePlace}
  \Delta Board \\
  w? : WORKER
\where
  w? \notin \dom loading \\
  began < 3 \\
  began' = began + 1 \\
  loading' = loading \cup \{ w? \mapsto began' \} \\
  held' = held \\
  shown' = shown \\
  staged' = staged \\
  everheld' = everheld \\
  evershown' = evershown \\
  offered' = offered \\
  overtaken' = overtaken
\end{schema}

\texttt{StoreBoard} is \texttt{BoardSlot.store}: the load has returned and its
board is offered to the slot, which keeps whichever of the two holds the later
place. The whole read-compare-write is one step, because it runs under the
slot's lock and contains nothing that can block --- a field read, a comparison
of two integers, and a field write.

The two disjuncts are \texttt{newer\_of}. If the incoming board began later it
displaces the one held; otherwise it is refused, and \texttt{overtaken} records
it. Equality cannot arise: each load takes its own place.

\begin{schema}{StoreBoard}
  \Delta Board \\
  w? : WORKER
\where
  w? \in \dom loading \\
  ( loading~w? > held \land held' = loading~w? \\
  \quad~ \land overtaken' = overtaken ) \lor \\
  \quad~ ( loading~w? \leq held \land held' = held \\
  \quad~ \land overtaken' = overtaken \cup \{ loading~w? \} ) \\
  ( held' > everheld \land everheld' = held' ) \lor \\
  \quad~ ( held' \leq everheld \land everheld' = everheld ) \\
  loading' = \{ w? \} \ndres loading \\
  began' = began \\
  shown' = shown \\
  staged' = staged \\
  evershown' = evershown \\
  offered' = offered
\end{schema}

% ============================================================
\section{Pushing}
\label{sec:push}
% ============================================================

A push writes the glass. It is two steps, and the design is entirely a
statement about what may happen between them.

\texttt{PushTake} settles what this push will write. It reads the slot and takes
the board the slot is holding \emph{now} --- not a board captured earlier, and
not the board this worker happens to have just loaded. Its guard, $staged =
\emptyset$, is the push gate: at most one push is settled-but-not-landed at any
time, which is what it means for a lock to be held across the write. That is the
whole of the fix, and it is two claims at once: the glass is written by one
pusher at a time, and each pusher reads the slot inside the region it writes in.

\begin{schema}{PushTake}
  \Delta Board \\
  w? : WORKER
\where
  staged = \emptyset \\
  staged' = \{ w? \mapsto held \} \\
  ( held > offered \land offered' = held ) \lor \\
  \quad~ ( held \leq offered \land offered' = offered ) \\
  began' = began \\
  loading' = loading \\
  held' = held \\
  shown' = shown \\
  everheld' = everheld \\
  evershown' = evershown \\
  overtaken' = overtaken
\end{schema}

There is one \texttt{PushTake} and not three, because under this design the
three pushes in the code are one act. A click's answer pushes what the slot
holds; a refresh pushes what the slot holds once its own board has been stored;
and a click with no board pushes what \texttt{NoBoard} renders, which is the
slot holding place $0$. Each is ``render what the slot is holding'', and the
model has no reason to tell them apart. Section~\ref{sec:oblig} says what the
implementation must do to make that true of the failure message as well.

\texttt{PushCommit} is the write landing at the display.

\begin{schema}{PushCommit}
  \Delta Board \\
  w? : WORKER
\where
  w? \in \dom staged \\
  shown' = staged~w? \\
  ( staged~w? > evershown \land evershown' = staged~w? ) \lor \\
  \quad~ ( staged~w? \leq evershown \land evershown' = evershown ) \\
  staged' = \{ w? \} \ndres staged \\
  began' = began \\
  loading' = loading \\
  held' = held \\
  everheld' = everheld \\
  offered' = offered \\
  overtaken' = overtaken
\end{schema}

% ============================================================
\section{Verification}
\label{sec:verify}
% ============================================================

Five safety properties. Each is checked by asking ProB whether its negation is
\emph{reachable}; a goal reported not found over the whole reachable state space
is the property proved.

\begin{description}
  \item[I1 --- the slot never goes backwards.] $held = everheld$. The board the
    slot holds is always the newest it has ever held. This re-verifies the
    shipped \texttt{BoardSlot} design and, more usefully, verifies that no
    operation added here writes the slot behind its own back.
  \item[I2 --- the glass never goes backwards.] $shown = evershown$. The display
    never returns to a board older than one it has already rendered. This is the
    property the answer's captured push violates.
  \item[I3 --- at rest, the glass shows the newest board anybody offered.]
    $staged = \emptyset \implies shown \geq offered$. Once no push is in flight,
    every push has either landed its board or found the glass already showing
    one at least as new. This is what rules out the degenerate reading of I2,
    under which a design that never pushed anything at all would pass.

    The antecedent is necessary and is not a weakening. Between settling on a
    board and landing it, a push has raised \texttt{offered} and not yet moved
    \texttt{shown}, so $shown < offered$ holds for the duration of every push
    ever made --- that inequality \emph{is} the push, not a defect. ProB reports
    the unconditioned form violated in four steps, and the trace is a
    \texttt{PushTake} with no \texttt{PushCommit} after it.
  \item[I4 --- the glass never runs ahead of the slot.] $shown \leq held$. The
    display never shows a board the slot does not hold. This is the property the
    refresh's push-before-store violates.
  \item[I5 --- the glass never shows a refused board.]
    $shown \notin overtaken$. The display never renders a board the slot
    examined and rejected as older than the one it was keeping.
\end{description}

\begin{verbatim}
  probcli docs/board_ordering.tex -model_check \
    -p DEFAULT_SETSIZE 2 -p MAXINT 4 -p TIME_OUT 60000 \
    -goal "held /= everheld"        # I1: must report goal NOT found
  probcli ... -goal "shown /= evershown"                # I2
  probcli ... -goal "staged = {} & shown < offered"     # I3
  probcli ... -goal "shown > held"                      # I4
  probcli ... -goal "shown : overtaken"                 # I5
\end{verbatim}

Against this model all five report the goal not found over 270 reachable states
and 657 transitions, exhaustively --- ``all open states visited'', every
operation covered.

The bare \texttt{-model\_check} additionally confirms that every operation is
covered and that the structural state predicate is preserved on every reachable
state.

\paragraph{Deadlock.}
There is no reachable state with nothing enabled. \texttt{PushTake} is enabled
whenever $staged = \emptyset$ and \texttt{PushCommit} whenever it is not, so one
of the two is always available and the model cannot wedge.

That is the state-machine result. The lock result is separate and is argued from
the code, because two of the three locks on this path are not modelled here. The
applet holds three:

\begin{itemize}
  \item \textbf{F}, \texttt{SingleFlight}'s. Every acquisition is
    \texttt{acquire(blocking=False)}. A lock that is never waited on cannot be
    the waiting edge of a cycle, so it cannot participate in a deadlock at all.
  \item \textbf{S}, \texttt{BoardSlot}'s. Its critical section is a field read,
    an integer comparison, and a field write. Nothing inside it blocks and
    nothing inside it acquires another lock, so S cannot be the \emph{held} edge
    of a cycle either.
  \item \textbf{P}, the push gate this design adds. It is held across the socket
    write, and --- because \texttt{PushTake} reads the slot inside it --- across
    an acquisition of S.
\end{itemize}

So exactly one nesting exists, $P \rightarrow S$, and the reverse never occurs:
no holder of S ever asks for P. One direction is not a cycle, so the design is
deadlock-free by acquisition order, and the model confirms the state machine
above it cannot stall.

This does change a claim the code currently makes. \texttt{board\_slot.py}'s
module docstring says of the slot's lock that ``there is no second lock to order
it against, so no acquisition order exists to get wrong''. After this change
there is a second lock and there is an order. The docstring must be corrected to
state it: S is the inner lock, P the outer, and nothing taken while holding S.

% ============================================================
\section{What the fix is, and what it is not}
\label{sec:oblig}
% ============================================================

Two candidate fixes were proposed for this defect. Neither closes it, and each
fails for its own reason; the model reaches a violating trace for both
(Section~\ref{sec:fidelity}).

\paragraph{A monotone push counter does not order the writes.}
The first candidate has the applet track the highest place it has pushed and
skip any push at a place no newer. The counter is claimed under a lock; the
socket write happens after the lock is released, because it must. So two pushers
can claim in the order $1, 2$ and land in the order $2, 1$, and the glass ends
at $1$. The counter records that a newer board was pushed while the display
shows an older one --- it makes the defect harder to see, not absent. Claiming a
place is not the same act as writing the glass, and only the second one matters.

\paragraph{Re-reading the slot does not close the gap either.}
The second candidate has the answer read the slot at push time instead of
pushing a board captured earlier. That shortens the window between the read and
the write but does not remove it, and the write is still unordered against other
pushers. A push that reads place $1$, is overtaken by a push that reads and
lands place $2$, and then lands its own $1$ regresses the glass exactly as
before.

\paragraph{What does close it.}
Re-reading the slot \emph{inside a region that is held across the write}. The
two candidates are the two halves of one fix, and each is inert without the
other: serialising the writes without re-reading still lands captured boards in
order but lands stale ones; re-reading without serialising reads the right board
and races to write it.

Given both, a third result follows, and it is the one worth acting on: \emph{the
place comparison is dead code}. I4 says the glass never runs ahead of the slot,
and I1 says the slot never goes back, so a board read from the slot inside the
push region is never older than what the glass is showing. There is no state in
which the comparison can refuse a push. A monotone \texttt{pushed} counter
beside the slot would therefore be a field that is written, read, and never once
acted on. Do not implement it.

\paragraph{The obligation this places on the implementation.}
The theorem holds for pushes whose content is read from the slot, which is why
the model has one \texttt{PushTake}. Three obligations make that true of the
code:

\begin{enumerate}
  \item \textbf{Every push reads the slot inside the push region.} The answer
    must not push a \texttt{HeldBoard} captured before the raise; the refresh
    must not push the board it loaded.
  \item \textbf{A refresh stores before it pushes.} Store first, then push what
    the slot holds. Push-first is what I4 forbids, and a refresh that pushed
    before storing would also push the \emph{previous} board under this design
    --- the user would wait out the query and never see its result.
  \item \textbf{The failure message must come from the slot.} The red message
    naming why \texttt{bd} could not be read is today the one push whose content
    the slot does not hold, and content that does not come from the slot escapes
    the theorem. The clean shape is for the state holding no board to carry the
    reason it holds none, so that ``render what the slot holds'' covers the
    failure as it covers the placeholder. If instead the message is pushed
    directly, it must be pushed inside the push region and only when
    $shown = 0$; I4 is what makes that guard sufficient.
\end{enumerate}

\paragraph{A boundary the fix does not close.}
The slot's lock is not held across the push --- deliberately, since the push is a
socket round trip. So a store may land between a push's read of the slot and its
write to the glass, leaving the glass one generation behind the slot: the user
clicked, and a warm-up that landed a millisecond later is held but not shown.
This state is reachable in the committed model and violates none of I1--I5,
because it is not a regression, not a refused board, and not a board the slot
does not hold --- it is ordinary staleness, indistinguishable from a warm-up
that landed a millisecond after the push instead of before it. Closing it would
mean either holding the slot's lock across a socket write, or having the pusher
re-check the slot after the write and push again when it moved. Neither is worth
its cost, and the boundary is recorded here rather than left to be rediscovered.

% ============================================================
\section{Fidelity: each variant reaches its own bad state}
\label{sec:fidelity}
% ============================================================

A model that cannot reproduce the defect it guards against is too abstract to
trust. Three variants do so, each in its own file, each differing from this one
only in Section~\ref{sec:push}.

\paragraph{The code as it stands ---
\texttt{docs/board\_ordering\_current\_buggy.tex}.}
Three staging operations, none gated: \texttt{AnswerStage} captures the board
the slot holds at the moment the click is acknowledged, \texttt{RefreshStage}
stages the board this worker loaded before the slot has seen it, and
\texttt{BlankStage} stages place $0$ whenever the slot holds no board. Over 3275
reachable states it reaches I2, I3, I4 and I5, and not I1. The I2 witness is the
one review found: a click captures place $1$; a refresh stores and lands place
$2$; the click lands its captured $1$ over it, and stands down. I4 needs no race
at all --- the refresh pushes before the store on every single refresh, so the
display is ahead of the slot for the duration of each one. I5 is a load that
began first, returned last, landed its board, and was then refused. That
document checks each witness against \texttt{SingleFlight} by hand, since this
family of models deliberately admits more interleavings than the code can
produce.

\paragraph{A monotone push counter alone ---
\texttt{docs/board\_ordering\_gate\_only\_buggy.tex}.}
Adds a \texttt{pushed} component and skips any staging at a place no newer,
leaving the writes unserialised. It reaches \emph{all four}: I2, I3, I4 and I5.
The counter closes nothing. Its I2 witness is the two-act argument in eight
steps --- claims at $1$ then $3$, writes at $3$ then $1$, and a counter reading
$3$ over a display showing $1$. ProB also reports \texttt{BlankStage} never
covered, because a strict counter over places can never admit place $0$: the
design would suppress the placeholder and the failure message outright,
including the first placeholder of a session.

\paragraph{Re-reading the slot alone ---
\texttt{docs/board\_ordering\_reread\_only\_buggy.tex}.}
The committed model with the single guard $staged = \emptyset$ deleted from
\texttt{PushTake}, and nothing else changed. Over 945 states it reaches I2 and
I3 --- the read is fresher, the write is still a race --- but not I4, I5 or I1.
That split is the evidence for the shape of the fix: re-reading the slot is what
keeps the glass from running ahead of the slot and from ever showing a refused
board, and serialising the writes is what keeps it from going backwards. One
guard separates a design that holds from one that does not.

\end{document}
